\FloatBarrier
\section{Discussion}
\label{sec:discussion}

\subsection{Principal Findings, Lesion Fidelity and the Cost of the Regulariser}
\label{sec:d-principal}
CT-supervised bone suppression can be made lesion-faithful, and its images help prediction in a specific way. Lesion absorption fell
steadily with the twin weight (Table~\ref{tab:twin}); on an audit set that no model had seen during development, the twin ensemble absorbed
a third of what the released CT-supervised model absorbed and less than a third of what the matched control absorbed
(Tables~\ref{tab:audit_means} and~\ref{tab:audit}), at non-inferior bone-region fidelity; and the effect replicated with two other encoders
(Table~\ref{tab:encoders}) and on paediatric radiographs (Section~\ref{sec:res-bdcxr}). Added as extra channels, the suppressed images
raised nodule detection on Node21, detection on VinBigData and classification on NIH (Table~\ref{tab:downstream}), and in the fusion study
they favoured focal findings, both in AUROC and in where the attribution maps pointed (Tables~\ref{tab:groups} and~\ref{tab:pointing}).

Several planned results did not hold. Our soft-tissue image alone did not improve any downstream task (Table~\ref{tab:downstream}); a
multi-stream classifier that read the radiograph, soft-tissue and bone images did not improve on the radiograph alone (H6,
Table~\ref{tab:fusion}); SSIM fell short of the reference model by more than its margin (Table~\ref{tab:fidelity}); real-nodule contrast was
higher than with the reference model on JSRT and lower on Node21 (Table~\ref{tab:cnr}); and conformal intervals calibrated on synthetic
data covered 14\,\% instead of 90\,\% of real bone pixels, 89\,\% after an exploratory recalibration on 24 real pairs
(Table~\ref{tab:reliability}).

\emph{What the audit adds.} Fidelity metrics and lesion absorption rank suppressors differently (Fig.~\ref{fig:pareto}). The v1 model has
the highest bone-region PSNR of the CT-supervised models and the highest absorption; ours and the reference model are within 0.06\,dB of
each other and a factor of three apart in absorption. BS-LDM, trained on dual-energy pairs from the evaluation hospital, is about 7\,dB ahead
in fidelity on that hospital's images, yet absorbed 0.42 of an inserted lesion in the development audit (at unpaired sites, so only
descriptive): fidelity to a dual-energy target does not protect what the target never showed. Part of that lead is the appearance of the
site's dual-energy images rather than suppression: a smooth calibration learned from 24 pairs of the site raised every CT-supervised model
by 2.1 to 2.4\,dB in the bone region without raising lesion absorption, and BS-LDM remained about 5\,dB ahead
(Section~\ref{sec:res-sitecal}; exploratory). Absorption must still be read beside suppression strength. The reference model removes more
rib-edge energy than ours (48\,\% against 30\,\%) and more non-rib energy (22\,\% against 8\,\%), so part of its higher absorption may
reflect stronger suppression in general.

\emph{What the regulariser costs.} On real paired radiographs the twin term costs no measurable bone-region fidelity up to $w=16$ (21.22 to
21.28\,dB across $w=0\ldots16$); at $w=32$ and 64 PSNR drops by about 0.3\,dB, so $w=16$, chosen on development data before the final
evaluation, is the usable operating point. The cost appears where the truth is exact: on the Qure.ai test split the bone error of the twin
ensemble is 41\,\% higher than that of the control (MAE 0.0041 against 0.0029; Supplementary Table~S5). The regulariser thus trades
synthetic bone accuracy for lesion invariance, a trade not visible on real images. The ablation shows that the effect needs gradients
through both twin branches and a consistency term confined to the pasted lesion; showing lesions during training without the twin term
(CLIA) did not lower absorption. The reference model is a SwinV2-B at $1024$~px trained with a far larger budget, so only the matched
no-twin control, which absorbs slightly more than the reference model (0.178 against 0.149), isolates the regulariser. After the
exploratory site calibration the no-twin control was slightly ahead of ours (23.74 against 23.58\,dB), a cost of the regulariser that the
uncalibrated comparison did not show; it was not tested. Diffusion models trained on real dual-energy pairs
\cite{sun2026bsldm,sun2026debonedit} need paired data from the target site, sample slowly (about 2.5 minutes per image for BS-LDM on a T4
GPU) and were not designed for lesion invariance.

\subsection{Why Suppressed Inputs Help Focal Findings}
\label{sec:d-focal}
The suppressed images are deterministic functions of the radiograph (Eq.~\eqref{eq:info}): the suppressor adds a transformation learned
from CT but no information, and a classifier trained on thousands of radiographs may learn whatever bone-invariant features its labels
reward. The results fit a modest, cheap prior rather than new information. Extra channels helped, whereas the soft image alone did not; the
multi-stream hybrid matched the radiograph alone on all three data sets, and on NIH its interval rules out a gain larger than 0.0024 in
macro-AUROC, so the null result is informative. The hybrid also matched the capacity control, and plain channel stacking did at least as
well as separate streams at a third of the encoder computation. Nodules and masses had the largest estimates on both adult data sets
($+0.008$ to $+0.009$ for the soft-tissue input), and the pointing game agreed: maps of the radiograph stream found focal lesions less often
than other findings, whereas maps of the soft-tissue stream found both about equally often. Not every signal favoured ours: on TBX11K the
reference model's soft image helped more, and on Node21 it raised real-nodule contrast more, so our less strongly suppressing model may
trade some contrast for lesion fidelity.

Three cautions apply. The NIH test split had already been used by the downstream classifiers of Study~I, so the NIH result of the fusion
study is a re-test with new models; VinBigData, the independent test, has a single merged nodule/mass class whose interval included zero;
and the detection gains are established for CPM, whereas the mAP@50 intervals include zero. The loss for diffuse findings appeared on NIH
only, mostly for infiltration; with report-derived NIH labels \cite{wang2017chestxray} and no such loss on the radiologist-labelled
VinBigData, these data cannot tell whether suppression removes cues for diffuse disease or interacts with label noise.

\subsection{Audit Circularity, Uncertainty and Clinical Translation}
\label{sec:d-translation}
\emph{Circularity.} The regulariser pastes lesions to teach invariance and the audit pastes lesions to measure it. The design limits the
circularity (Section~\ref{sec:m-audit}): audit lesions come from held-out LIDC-IDRI test CTs, audit images are real radiographs, the
confirmatory audit uses rib masks that no model has seen, and evidence that does not use insertion points the same way, namely the contrast
of real JSRT nodules (ours above the reference model and the no-twin control) and the Node21 detection gain of the three-channel input. These
choices limit the circularity but cannot remove it: a model that recognises pasted signal in general could still score well.

\emph{Uncertainty.} The bone-image intervals held where they were calibrated and failed under the shift to real radiographs; 24 real pairs
from the target site sufficed to restore coverage here, although transfer between sites is untested. For the classifiers, conformal risk
control kept its guarantee for every input, but label sets holding about half of the 14 NIH labels are of little use to a reader, and the
suppressor's own uncertainty made abstention worse in every configuration: uncertainty about the bone image is not uncertainty about the
diagnosis.

\emph{Recommendations.} A bone suppressor should be audited for lesion absorption, reported next to a suppression measure, before it is
used; suppressed images should be shown beside the original rather than instead of it; and the ensemble spread should serve as a warning
flag rather than a calibrated interval unless recalibrated on real pairs (it rose 2.7-fold on paediatric radiographs). Detection and
classification models should receive suppressed images as extra channels next to the radiograph rather than as separate streams, should
not be expected to gain more than small amounts of overall accuracy, and should be tested per finding against a control of the same
capacity. Fidelity comparisons between CT-supervised and dual-energy-trained suppressors should separate the appearance of the test site
from suppression, for example with a calibration on a few development pairs reported beside the uncorrected numbers.

\subsection{Limitations and Future Research}
\label{sec:limitations}
Each limitation is stated with its consequence for how the results must be read.
\begin{enumerate}
\item \emph{Synthetic supervision and one source of real pairs.} The bone image of a real radiograph is never observed. Real-image fidelity
 is measured against dual-energy soft-tissue images from one hospital and one acquisition set-up (GE Discovery XR656), which differ from
 the radiograph in tone and fine texture, so raw and band-limited metrics are both reported and calibrated metrics are not absolute.
\item \emph{The audit is linear and additive, and shares its mechanism with the regulariser.} Lesions are added in image units; beam
 hardening, scatter and tone mapping are not modelled, so $A$ measures the leak of an added signal rather than clinical detectability, and
 it is read next to a suppression gain because $A=0$ is reached by doing nothing. The real-nodule contrast and downstream subgroup analyses
 limit, but do not remove, the circularity of inserting lesions both in training and in the audit.
\item \emph{Development use of SZCH.} All 75 SZCH images were used in development (audit) and 24 for paired decisions, so only the 51
 held-out pairs support fidelity claims; the recipe was revised once after the SZCH pairs had been looked at. The exploratory site
 calibration was designed after H2$'$ was known, re-uses the 51 held-out pairs and is specific to one hospital and scanner, whose own 24
 development pairs it was fitted on.
\item \emph{Resolution, backbone and operating point.} Our suppressors predict at 512\,px with small encoders, the reference model at
 1024\,px with SwinV2-B. The SSIM gap sits in fine texture (with a band limit of $\sigma=4$ it shrinks from 0.021 to 0.002); whether
 prediction at 1024\,px closes it was not tested. Only $w=16$ was evaluated with the full protocol.
\item \emph{Uncertainty.} Conformal calibration treats pixels as exchangeable and was calibrated on synthetic data; the recalibration on
 real pairs was chosen after the shift was seen and used 24 pairs from one hospital. Ensemble members share training data.
\item \emph{Scope and external validity.} Training and paired evaluation use adult frontal radiographs; BDCXR-3257 is a single-centre
 paediatric set with image-level labels and no patient identifiers, and its audit inserts adult nodules into children's radiographs. There
 is no lateral imaging and no reader study. On JSRT the bone overlap of each audit site came from the reference model's bone map.
 Downstream models were not tuned per input, the overlap analysis is not estimable on TBX11K, and the downstream gains are small and were
 measured with one classifier and one detector. The models are for research use only.
\item \emph{Fusion study.} One encoder (DenseNet-121) and one resolution ($256\times256$\,px) were tested; heavier encoders or full
 resolution may use the extra streams differently. The NIH test split served the downstream experiments of Study~I as well, its labels are
 report-derived, and VinBigData has one focal class. The per-class DeLong tests treat NIH images as independent although patients
 contribute several images, and the pointing game depends on the attribution method, the 15\,px tolerance and the eight annotated NIH
 classes.
\end{enumerate}

\emph{Future research.} Future work should test prediction at $1024$~px, recalibration on real pairs from more sites, other lesion types,
detectors that receive the suppressed images as extra channels at full resolution, a reader study and a prospective evaluation.
